\ncdmodultest{3}{Fællesskab, Traditioner og Sprog}{全真模考卷三：志愿服务、圣诞午餐与丹麦语学习}{社区志愿者招募 / 烧烤派对上的不速之客 / 丹麦圣诞午餐文化 / 三位外国人谈如何学好丹麦语 / 新运动爱好回信与网购维权 / 大学学习与冬泳口语及酒店实习信息差}{反身物主代词 sin/sit/sine vs deres · 相互代词 hinanden vs sig · 目的 for at vs 让步 selvom · 间接引语 at/om/hv- · 对机构用复数 I/jer}

% =====================================================================
% Del 1: Læsning (40 minutter · 24 point)
% =====================================================================
\ncdsubsec{Del 1: Læsning（阅读理解 · 40 minutter · 24 point）}

\begin{tipbox}{考场规则与限时建议（Instruktion til læseprøven）}
\small
本部分共四道大题，总计 \textbf{40 分钟}（\textbf{禁用任何词典或手机}）。满分 \textbf{24 分}，\textbf{16--24 分及格（Bestået）}；0--15 分不及格（Ikke bestået）。建议考场配速：\textbf{Opgave 1 (8 min) $\rightarrow$ Opgave 2 (8 min) $\rightarrow$ Opgave 3 (8 min) $\rightarrow$ Opgave 4 (12 min)}，留 4 分钟复查。
\end{tipbox}

% ---------------------------------------------------------------------
\subsection*{Opgave 1 (NY · 6 point): \textit{Bliv frivillig i lokalområdet}}

\begin{grammatikbox}{Instruktion — Læsning, Opgave 1}
\small
\begin{itemize}
  \item Læs de seks små tekster om personerne, der gerne vil lave frivilligt arbejde (1--6).
  \item Find det opslag / den annonce (B--I), der passer til hver tekst.
  \item Skriv bogstavet på annoncen, der passer til teksten. Der er \textbf{to annoncer, du ikke skal bruge}. Se eksemplet (0).
\end{itemize}
\end{grammatikbox}

\begin{enumerate}
  \item[0.] Astrid er gået på pension og elsker at læse romaner og historie. Hun savner kontakten til andre mennesker og vil gerne bruge et par eftermiddage om ugen i en hyggelig butik, hvor hun kan sortere bøger og hjælpe kunder, så overskuddet går til børn i nød. \hfill \textcolor{NCDRed}{\textbf{[Eksempel: Svar A]}}

  \item[1.] Emil arbejder på kontor om dagen, men han bor alene og har god tid om aftenen. Hans egen bedstemor boede langt væk, så han savner kontakten til ældre generationer. Han vil gerne mødes fast med én ældre person en gang om ugen for at drikke kaffe og snakke sammen. \hfill \textbf{Svar:}\ \rule{1.2cm}{0.4pt}

  \item[2.] Nadia studerer matematik på universitetet og er rigtig dygtig til tal. Før hun flyttede til Danmark, var hun hjælpelærer på en skole. Hun vil gerne hjælpe unge elever efter skoletid, så de bedre kan forstå deres hjemmearbejde, men hun har kun tid om tirsdagen. \hfill \textbf{Svar:}\ \rule{1.2cm}{0.4pt}

  \item[3.] Tobias bor i en lille lejlighed midt i byen, hvor man ikke må have husdyr. Han er vokset op på landet og savner især at gå lange ture med hunde. Derfor leder han efter et sted, hvor han kan hjælpe til med at passe og lufte dyr i weekenden. \hfill \textbf{Svar:}\ \rule{1.2cm}{0.4pt}

  \item[4.] Samira taler både dansk, arabisk og engelsk, og hun synes, det er spændende at møde mennesker fra hele verden. Hun vil gerne mødes med en nyankommen flygtning i lokalområdet for at hjælpe med at øve dansk gennem helt almindelige samtaler om job, kultur og hverdagsliv. \hfill \textbf{Svar:}\ \rule{1.2cm}{0.4pt}

  \item[5.] Mads bekymrer sig meget om naturen og dyrene i havet. Han sidder foran en computerskærm hele ugen, så i sin fritid vil han helst være ude i den friske luft ved vandet og gøre noget praktisk for at fjerne plastik og affald sammen med andre. \hfill \textbf{Svar:}\ \rule{1.2cm}{0.4pt}

  \item[6.] Katrine læser til sygeplejerske og har fri hver torsdag formiddag. Hun vil gerne bruge nogle timer på et sygehus, hvor hun kan byde patienter og pårørende velkommen ved indgangen og hjælpe dem med at finde den rigtige afdeling, så de føler sig trygge. \hfill \textbf{Svar:}\ \rule{1.2cm}{0.4pt}
\end{enumerate}

\vspace{0.2em}
{\small
\setlength{\LTleft}{\fill}
\setlength{\LTright}{\fill}
\renewcommand{\arraystretch}{1.35}
\begin{longtable}{@{} >{\raggedright\arraybackslash}p{8.2cm} @{\hspace{0.6cm}} >{\raggedright\arraybackslash}p{8.2cm} @{}}
\toprule
\multicolumn{2}{@{} >{\raggedright\arraybackslash}p{17.0cm} @{}}{%
  \badge{NCDRed}{A (Eksempel)}\ \textbf{\color{NCDRed}Børnenes Genbrugsboghandel -- Vi mangler frivillige!}\newline
  Elsker du bøger, og har du lyst til at hjælpe andre? I vores lille boghandel midt i byen sælger vi brugte romaner, børnebøger og fagbøger. Som frivillig står du for at sætte bøger på hylderne, tale med kunderne og passe kassen 3--4 timer om eftermiddagen. Alle penge går direkte til udsatte børn.%
} \\
\midrule
\endfirsthead
\toprule
\multicolumn{2}{@{} >{\raggedright\arraybackslash}p{17.0cm} @{}}{%
  \textbf{\color{NCDNavy}Opslag / Annoncer B--I（续表）}%
} \\
\midrule
\endhead
\bottomrule
\endfoot
\badge{NCDNavy}{B}\ \textbf{\color{NCDNavy}Frivillig hundelufter hos Dyreinternatet Vestermark}\newline
Hos Dyreinternatet Vestermark tager vi os af hunde og katte, der venter på et nyt hjem. Vi søger pålidelige frivillige over 18 år, som har lyst til at komme lørdag eller søndag formiddag. Din opgave er at gå en lang tur med vores hunde i skoven, give dem frisk vand og lege med dem, så de får masser af motion og tryghed.
&
\badge{NCDNavy}{F}\ \textbf{\color{NCDNavy}Patientguider på Regionshospitalet}\newline
Det kan være svært og forvirrende at finde vej på et stort hospital, når man er syg eller nervøs. Derfor søger vi venlige frivillige til vores team ved hovedindgangen mandag til fredag mellem kl. 8.30 og 13.00. Du skal tage godt imod patienter og gæster, vise vej til den rigtige afdeling eller følge dem hen til venteværelset. \\
\midrule
\badge{NCDNavy}{C}\ \textbf{\color{NCDNavy}Fælleshaven Grøn Oase}\newline
Kan du lide at arbejde med jorden og se ting vokse? I vores lokale fælleshave midt i kvarteret dyrker vi økologiske gulerødder, tomater og blomster. Vi mødes hver lørdag formiddag for at vande planterne, luge ukrudt og bygge nye trækasser sammen. Du behøver ikke have grønne fingre i forvejen -- vi lærer hinanden op og spiser frokost sammen bagefter.
&
\badge{NCDNavy}{G}\ \textbf{\color{NCDNavy}Frivillige søges til Ungdommens Lektiecafé}\newline
Har du lyst til at hjælpe børn og unge i alderen 10--16 år med at blive bedre i skolen? Hver tirsdag mellem kl. 15.00 og 17.30 åbner vi vores lektiecafé på biblioteket. Her laver vi sjove opgaver sammen, og vi har især brug for nogen, der kan forklare matematik, fysik eller dansk. Du behøver ikke være uddannet lærer, bare tålmodig og venlig. \\
\midrule
\badge{NCDNavy}{D}\ \textbf{\color{NCDNavy}Bliv besøgsven for en ældre borger}\newline
Mange ældre mennesker i vores kommune sidder meget alene og savner nogen at tale med. Som besøgsven bliver du matchet med en ældre dame eller herre i nærheden af dig. I mødes cirka to timer en eftermiddag eller aften hver uge over en kop kaffe, går en kort tur eller spiller kort. Det koster ikke noget, men det giver stor glæde for jer begge.
&
\badge{NCDNavy}{H}\ \textbf{\color{NCDNavy}IT-hjælp på Kulturhuset}\newline
Har du styr på telefoner, apps og MitID, og har du lyst til at hjælpe seniorer? Hver onsdag fra kl. 10.00 til 12.00 holder vi åben IT-café i Kulturhuset, hvor ældre borgere kan komme med spørgsmål om digital post, netbank eller beskeder på telefonen. Du sidder ved et bord sammen med andre frivillige og forklarer tingene roligt trin for trin. \\
\midrule
\badge{NCDNavy}{E}\ \textbf{\color{NCDNavy}Bliv sprogmakker for en flygtning}\newline
Når man kommer til Danmark som flygtning, kan det være svært at få talt nok dansk uden for sprogskolen. I vores forening matcher vi en lokal frivillig med en ny borger. I aftaler selv, hvornår og hvor I mødes -- for eksempel på biblioteket eller til en gåtur -- og snakker sammen på dansk om arbejde, traditioner og livet i byen.
&
\badge{NCDNavy}{I}\ \textbf{\color{NCDNavy}Ren Strand -- Kom med ud til kysten!}\newline
Hver anden søndag mødes en stor gruppe frivillige på stranden ved Nørrebugt for at samle plastik, dåser og andet affald ind, før det ender i havet. Vi sørger for handsker, poser og varm kaffe efter arbejdet. Det er gratis at være med, og alle er velkomne -- uanset alder. Sammen gør vi naturen renere for både fugle, fisk og badegæster. \\
\end{longtable}
}

% ---------------------------------------------------------------------
\subsection*{Opgave 2 (5 point): \textit{En uventet gæst til grillfesten}}

\begin{grammatikbox}{Instruktion — Læsning, Opgave 2}
\small
Læs teksten \textit{En uventet gæst til grillfesten}. I hvert afsnit er der \textbf{én sætning, der ikke passer i afsnittet}. Find den sætning, der ikke passer. Streg sætningen ud. Se eksemplet i afsnit 0.
\end{grammatikbox}

\textbf{(0)} Lørdag den 14. juli klokken 17.30 havde Henrik og hans kone, Louise, inviteret seks gode venner til sommerens første grillfest i deres have. Solen skinnede fra en blå himmel, og alle gæsterne sad allerede ude på terrassen med kolde drikkevarer. \textcolor{NCDRed}{\textbf{[Eksempel — streges ud:]}}\ \textit{Det sneede kraftigt hele aftenen.} Henrik stod ved grillen og stegte pølser og kylling, mens duften af sommermad spredte sig ud over hækken.

\textbf{(1)} Klokken halv syv var maden næsten klar, og Louise gik ind i køkkenet for at hente salatskålen og det nybagte brød. Pludselig hørte gæsterne en mærkelig lyd bag de høje buske ved havelågen. Alle lagde deres glas fra sig og kiggede nysgerrigt over mod hjørnet af græsplænen. Ingen lagde mærke til den høje lyd. I det samme hoppede en lille brun hund med et rødt halsbånd ud fra buskene og løb direkte hen mod grillen.

\textbf{(2)} Hunden var tydeligvis meget sulten, for den snusede ivrigt rundt om Henriks ben og prøvede at nå op til pølserne på grillen. Henrik løftede forsigtigt den lille hund op, så den ikke brændte sig på det varme metal. Den logrede glad med halen og slikkede Henrik på hånden. Den store sorte hund gøede surt og bed alle gæsterne i benene. Louise kom hurtigt ud fra køkkenet med en skål koldt vand til den uventede gæst.

\textbf{(3)} Louise kiggede nærmere på hundens røde halsbånd og fandt et lille metalskilt med et navn og et telefonnummer. Der stod, at hunden hed Buster, og at den boede på Bakkevej 12 lige bag ved deres have. Henrik fandt straks sin telefon frem og ringede til nummeret på skiltet. Halsbåndet havde hverken navn eller adresse på. En ung mand ved navn Mikkel tog telefonen og lød meget lettet, da Henrik fortalte, at Buster sad trygt på terrassen.

\textbf{(4)} Fem minutter senere ringede det på hoveddøren, og Mikkel stod udenfor med en lang hundesnor i hånden. Han forklarede med et smil, at Buster havde gravet et lille hul under hækken, mens Mikkel vandede blomster i forhaven. Henrik viste ham ud på terrassen, hvor Buster straks løb hen til sin ejer. Mikkel blev sur og nægtede at tage hunden med hjem. Alle gæsterne lo venligt af historien, og Louise tilbød Mikkel et glas kold limonade, før han gik hjem.

\textbf{(5)} Da Mikkel og Buster var gået tilbage til Bakkevej, satte alle vennerne sig endelig til bords på terrassen. Kyllingen og salaten smagte fantastisk, og de talte længe om deres sjove besøg midt under aftensmaden. Stemningen var varm og afslappet, og ingen havde travlt med at gå hjem. Derfor aflyste Henrik festen og sendte alle vennerne hjem uden aftensmad. Først ved ellevetiden sagde gæsterne tak for en dejlig aften og cyklede hjem i den lyse sommernat.

% ---------------------------------------------------------------------
\subsection*{Opgave 3 (6 point): \textit{Julefrokost -- en vigtig dansk tradition}}

\begin{grammatikbox}{Instruktion — Læsning, Opgave 3}
\small
Læs teksten \textit{Julefrokost -- en vigtig dansk tradition}. Skriv de ord, der mangler. Du skal kun bruge hvert ord én gang. Der er \textbf{fire ord, du ikke skal bruge}. Se eksemplet i første linje.
\end{grammatikbox}

\begin{emailbox}{Julefrokost -- en vigtig dansk tradition}
I november og december holder mange danskere en særlig fest, som \underline{\textbf{hedder}} julefrokost. Man holder både julefrokost med sin familie, sine venner (1)\ \rule{2.1cm}{0.4pt} sine kolleger på arbejdet.

\vspace{0.3em}
Til en traditionel julefrokost står der mange (2)\ \rule{2.1cm}{0.4pt} på bordet, for eksempel sild, fiskefilet, frikadeller og flæskesteg. Mange drikker også juleøl eller snaps til maden, selvom nogle hellere vil have vand eller sodavand.

\vspace{0.3em}
Efter maden leger man tit pakkeleg, hvor alle gæsterne har taget en lille (3)\ \rule{2.1cm}{0.4pt} med. Man slår med en terning for at vinde, og stemningen bliver hurtigt både (4)\ \rule{2.1cm}{0.4pt} og højlydt.

\vspace{0.3em}
En firmajulefrokost er populær, fordi kollegerne mødes uden for kontoret i afslappet tøj. Det er en god måde at lære (5)\ \rule{2.1cm}{0.4pt} bedre at kende før det nye år. Hvis man fester sent, skal man (6)\ \rule{2.1cm}{0.4pt} huske at tage toget eller en taxa hjem i stedet for at køre i bil.
\end{emailbox}

\begin{tipbox}{Ordboks（共 11 个候选词：1 个例词 + 6 个正解 + 4 个干扰词）}
\centering
\textbf{\color{NCDNavy}retter \quad · \quad når \quad · \quad hedder \quad · \quad altid \quad · \quad men \quad · \quad gave \quad · \quad hinanden \quad · \quad havde \quad · \quad sjov \quad · \quad mellem \quad · \quad og}
\end{tipbox}

% ---------------------------------------------------------------------
\subsection*{Opgave 4 (NY · 7 point): \textit{Tre udlændinge fortæller om at lære dansk}}

\begin{grammatikbox}{Instruktion — Læsning, Opgave 4}
\small
\begin{itemize}
  \item Læs de tre tekster, og læs spørgsmålene (1--7).
  \item Find den person, der passer til hvert af de syv spørgsmål.
  \item Sæt kryds (\textbf{X}) ud for den person, der passer til spørgsmålet. Se eksemplet (0).
\end{itemize}
\end{grammatikbox}

\begin{tcolorbox}[enhanced, breakable, colback=NCDBgGray, colframe=NCDNavy, boxrule=0.8pt, arc=3pt, left=8pt, right=8pt, top=7pt, bottom=7pt]
\small
\renewcommand{\arraystretch}{1.28}
\begin{tabularx}{\linewidth}{@{} >{\raggedright\arraybackslash}X >{\centering\arraybackslash}p{2.0cm} >{\centering\arraybackslash}p{2.0cm} >{\centering\arraybackslash}p{2.0cm} @{}}
\toprule
\textbf{Spørgsmål} & \textbf{A. Marco} & \textbf{B. Yuna} & \textbf{C. Elena} \\
\midrule
0. Hvem var i starten for nervøs til at tale dansk? & A & B\ \textcolor{NCDRed}{\textbf{X}} & C \\
\midrule
1. Hvem bor sammen med en dansker? & A \ \ \underline{\hspace{0.4cm}} & B \ \ \underline{\hspace{0.4cm}} & C \ \ \underline{\hspace{0.4cm}} \\
2. Hvem øver dansk i pauserne på sit arbejde? & A \ \ \underline{\hspace{0.4cm}} & B \ \ \underline{\hspace{0.4cm}} & C \ \ \underline{\hspace{0.4cm}} \\
3. Hvem bruger litteratur for børn til at lære nye ord? & A \ \ \underline{\hspace{0.4cm}} & B \ \ \underline{\hspace{0.4cm}} & C \ \ \underline{\hspace{0.4cm}} \\
4. Hvem øver sin skrivning hver eneste dag? & A \ \ \underline{\hspace{0.4cm}} & B \ \ \underline{\hspace{0.4cm}} & C \ \ \underline{\hspace{0.4cm}} \\
5. Hvem træner dansk ved at hjælpe ulønnet i en forretning? & A \ \ \underline{\hspace{0.4cm}} & B \ \ \underline{\hspace{0.4cm}} & C \ \ \underline{\hspace{0.4cm}} \\
6. Hvem ser fjernsyn med tekst på skærmen for at lære udtalen? & A \ \ \underline{\hspace{0.4cm}} & B \ \ \underline{\hspace{0.4cm}} & C \ \ \underline{\hspace{0.4cm}} \\
7. Hvem hører dansk tale, mens de laver husarbejde eller cykler? & A \ \ \underline{\hspace{0.4cm}} & B \ \ \underline{\hspace{0.4cm}} & C \ \ \underline{\hspace{0.4cm}} \\
\bottomrule
\end{tabularx}
\end{tcolorbox}

\textbf{\color{NCDNavy}A. Marco}\\
Jeg kom fra Italien til Aarhus for to år siden for at arbejde som ingeniør. I det første år lærte jeg næsten ikke noget dansk, selvom jeg gik på sprogskole to aftener om ugen. På mit arbejde talte alle nemlig engelsk til møderne. For et halvt år siden lavede jeg derfor en fast aftale med to danske kolleger: I frokostpausen og over kaffen snakker vi kun dansk med hinanden. De må ikke skifte til engelsk, selvom jeg taler langsomt. Om aftenen ser jeg også danske krimiserier på tv med danske undertekster. Det hjælper mig enormt meget med at forstå, hvordan de skrevne ord lyder i virkeligheden. Jeg laver stadig grammatikfejl, men nu tør jeg endelig tale dansk med mine naboer og i butikkerne.

\vspace{0.3em}
\textbf{\color{NCDNavy}B. Yuna}\\
Da jeg flyttede fra Sydkorea til Odense for at studere design, var jeg meget genert. Jeg var bange for at sige noget forkert, så jeg turde næsten ikke åbne munden på dansk. En dag foreslog min lærer, at jeg skulle finde et sted at øve mig uden for skolen. Nu arbejder jeg frivilligt fire timer hver lørdag i en lokal genbrugsbutik med tøj og møbler. Der kommer mange snakkesalige kunder, og gennem korte samtaler om priser og vejr er jeg blevet meget tryggere ved at tale. Derudover låner jeg tynde børnebøger på biblioteket og læser dem højt hjemme på mit værelse. Sætningerne er enkle, og de samme vigtige ord bliver gentaget mange gange, så jeg nemmere kan huske dem.

\vspace{0.3em}
\textbf{\color{NCDNavy}C. Elena}\\
Jeg flyttede fra Spanien til Roskilde for halvandet år siden sammen med min danske mand. Jeg arbejder hjemmefra som grafiker for spanske kunder, så jeg sidder alene ved computeren hele dagen og hører ikke meget dansk på jobbet. Derfor lytter jeg til dansk radio og lette podcasts i næsten to timer hver dag, mens jeg cykler, gør rent eller laver mad. Selvom jeg ikke forstår hvert eneste ord, har mit øre vænnet sig til sprogets rytme og de bløde endelser. Hver aften skriver jeg desuden fem korte sætninger i en notesbog om min dag. Bagefter læser min mand dem igennem og forklarer mig, hvordan ordstillingen skal være, så jeg ikke glemmer reglerne. Det giver mig meget mere ro, når jeg skal tale dansk i byen.

% =====================================================================
% Del 2: Skrivning (minimum 70 ord)
% =====================================================================
\ncdsubsec{Del 2: Skrivning（书面表达 · minimum 70 ord）}

\subsection*{Opgave 1（官方必考私人邮件回复 · minimum 70 ord）}

\textbf{Situation:} Du har fået en mail fra din ven, Jonas.\\
\textbf{Opgave:} Læs e-mailen, og skriv et svar til Jonas. Du skal svare på alle Jonas’ spørgsmål. Du skal skrive \textbf{minimum 70 ord}（建议 85--115 词）.

\begin{emailbox}{Indbakke — Fra: Jonas \quad | \quad Emne: Ny fritidsaktivitet?}
\textbf{Hej}

Jeg har hørt, at du er begyndt til en ny sport eller fritidsaktivitet! Hvad går du til? Hvor ofte og hvorhenne gør du det? Hvorfor valgte du lige det? Og hvordan har de andre på holdet taget imod dig?

\textbf{Mange hilsner}\\
\textbf{Jonas}
\end{emailbox}

\begin{tcolorbox}[enhanced, unbreakable, colback=white, colframe=NCDNavy, coltitle=white, fonttitle=\bfseries, title={Dit svar — Til: Jonas \quad | \quad Emne: Ny fritidsaktivitet}, arc=2pt, boxrule=0.8pt, left=10pt, right=10pt, top=5pt, bottom=5pt]
\rule{\linewidth}{0.4pt}\\[0.38em]
\rule{\linewidth}{0.4pt}\\[0.38em]
\rule{\linewidth}{0.4pt}\\[0.38em]
\rule{\linewidth}{0.4pt}
\end{tcolorbox}

\subsection*{Bonus Opgave 1B（半正式实用邮件拓展 · 网购维权与催件 · minimum 70 ord）}

\textbf{Situation:} For 12 dage siden bestilte du \textbf{en vinterjakke} i en dansk netbutik (ordrenummer: \textbf{33918}). Ifølge ordrebekræftelsen skulle pakken komme \textbf{i løbet af 2--4 hverdage}. I onsdags fik du en besked om, at ordren var pakket, men du har ikke hørt noget \textbf{siden i onsdags}. Du skal snart ud at rejse og har brug for jakken.\\
\textbf{Opgave:} Skriv en høflig, halvformel e-mail til \textbf{Kundeservice} på \textbf{minimum 70 ord}. Forklar problemet, spørg hvornår pakken bliver leveret, og husk at bruge \textbf{flertal om firmaet (\textit{I / jer / jeres})} hele vejen igennem!

\begin{tcolorbox}[enhanced, unbreakable, colback=white, colframe=NCDTeal, coltitle=white, fonttitle=\bfseries, title={Dit svar — Til: Kundeservice \quad | \quad Emne: Manglende levering -- ordrenummer 33918}, arc=2pt, boxrule=0.8pt, left=10pt, right=10pt, top=5pt, bottom=5pt]
\rule{\linewidth}{0.4pt}\\[0.38em]
\rule{\linewidth}{0.4pt}\\[0.38em]
\rule{\linewidth}{0.4pt}\\[0.38em]
\rule{\linewidth}{0.4pt}
\end{tcolorbox}

% =====================================================================
% Del 3: Mundtlig kommunikation (口语交际全真题卡)
% =====================================================================
\clearpage
\ncdsubsec{Del 3: Mundtlig kommunikation（口语交际全真题卡）}

\subsection*{Mundtlig Opgave 1 (Emneopgave med Mindmap)}

\begin{tipbox}{官方考规提示（Opgave 1）}
\small
考生须自带 \textbf{2 张 Mindmap 主题卡}进入考场（其中至少 1 张属于工作 / 实习 / 大学学习 / 语言学校主题）。卡片上\textbf{只能填写关键词或短语（stikord），严禁写完整句子}。考官现场抽签决定其中 1 张：\textbf{Del 1} 自主陈述 \textbf{1--2 分钟}；\textbf{Del 2} 考官根据你的陈述连续追问 \textbf{3--4 分钟}（两部分均须通过）。
\end{tipbox}

\begin{tcolorbox}[enhanced, unbreakable, colback=NCDBgBlue, colframe=NCDNavy, coltitle=white, fonttitle=\bfseries, title={Mindmap 3A：Mit studie（我的学业 / 大学学习 · Stikord 备考卡）}, arc=3pt, boxrule=0.8pt, left=10pt, right=10pt, top=5pt, bottom=5pt]
\small
\renewcommand{\arraystretch}{1.18}
\begin{tabularx}{\linewidth}{@{} >{\raggedright\arraybackslash}X @{\hspace{0.5cm}} >{\raggedright\arraybackslash}X @{}}
\toprule
\multicolumn{2}{@{} >{\centering\arraybackslash}p{15.8cm} @{}}{\textbf{\color{NCDNavy}中心主题：Mit studie}} \\
\midrule
\textbf{\color{NCDTeal}1. uddannelse / sted}\newline
· datalogi / it-uddannelse\newline
· Københavns Universitet (KU)\newline
· Nørrebro · 2. semester &
\textbf{\color{NCDTeal}2. dage / tid}\newline
· mandag--torsdag · kl. 9.00--15.00\newline
· fredag: selvstudie hjemme \\
\midrule
\textbf{\color{NCDTeal}3. fag / lærere}\newline
· programmering · matematik · projekter\newline
· dygtige og hjælpsomme lærere\newline
· afslappet tone &
\textbf{\color{NCDTeal}4. studiekammerater}\newline
· danskere og internationale\newline
· lille studiegruppe · hjælper hinanden \\
\midrule
\textbf{\color{NCDTeal}5. pauser}\newline
· frokost kl. 12.00 · kantinen\newline
· kaffe og snak med gruppen &
\textbf{\color{NCDTeal}6. transport til studie}\newline
· cykel · 20 minutter langs søerne\newline
· metro ved kraftig regn \\
\bottomrule
\end{tabularx}

\vspace{0.25em}
\textbf{\color{NCDRed}Del 2 考官追问模拟题（4 题）：}
\begin{enumerate}[itemsep=0.12em, topsep=0.15em]
  \item \textit{Du siger, at du har dygtige lærere og gruppeprojekter. Hvilket fag kan du bedst lide, og hvorfor?}
  \item \textit{Hvilket sprog taler du med dine studiekammerater i pauserne?}
  \item \textit{Hvad vil du gerne arbejde med, når du er færdig med din uddannelse?}
  \item \textit{Er det svært at have tid til både studie og dansk sprogskole?}
\end{enumerate}
\end{tcolorbox}

\begin{tcolorbox}[enhanced, unbreakable, colback=NCDBgBlue, colframe=NCDNavy, coltitle=white, fonttitle=\bfseries, title={Mindmap 3B：Min fritidsaktivitet: Vinterbadning（我的业余活动：冬泳 · Stikord 备考卡）}, arc=3pt, boxrule=0.8pt, left=10pt, right=10pt, top=5pt, bottom=5pt]
\small
\renewcommand{\arraystretch}{1.18}
\begin{tabularx}{\linewidth}{@{} >{\raggedright\arraybackslash}X @{\hspace{0.5cm}} >{\raggedright\arraybackslash}X @{}}
\toprule
\multicolumn{2}{@{} >{\centering\arraybackslash}p{15.8cm} @{}}{\textbf{\color{NCDNavy}中心主题：Min fritidsaktivitet: Vinterbadning}} \\
\midrule
\textbf{\color{NCDTeal}1. sted / by}\newline
· Vinterbadeklubben Isfuglen\newline
· Amager Strandpark · København\newline
· badebro og sauna &
\textbf{\color{NCDTeal}2. dage / tid}\newline
· tirsdag og torsdag morgen · kl. 7.00--7.45\newline
· ofte søndag formiddag \\
\midrule
\textbf{\color{NCDTeal}3. pris / gratis}\newline
· medlemskab · 650 kr. om året\newline
· inkl. omklædning og varm sauna &
\textbf{\color{NCDTeal}4. alene / sammen med}\newline
· sammen med to naboer\newline
· faste medlemmer · hyggeligt fællesskab \\
\midrule
\textbf{\color{NCDTeal}5. aktiviteter}\newline
· kort dukkert i koldt vand\newline
· 15 min. i sauna · varm te bagefter &
\textbf{\color{NCDTeal}6. transport til aktiviteten}\newline
· cykel · 12 minutter\newline
· badetøj og stort håndklæde i rygsækken \\
\bottomrule
\end{tabularx}

\vspace{0.25em}
\textbf{\color{NCDRed}Del 2 考官追问模拟题（4 题）：}
\begin{enumerate}[itemsep=0.12em, topsep=0.15em]
  \item \textit{Hvorfor begyndte du egentlig at vinterbade? Var det ikke meget koldt i starten?}
  \item \textit{Hvordan føles kroppen efter saunaen og det kolde vand om morgenen?}
  \item \textit{Hvad siger dine venner eller din familie til, at du bader om vinteren?}
  \item \textit{Laver I også noget andet socialt i klubben ud over at bade?}
\end{enumerate}
\end{tcolorbox}

\subsection*{Mundtlig Opgave 2 (Informationskløft · 满分 6 分，4--6 分及格)}

\begin{tipbox}{考场规则提醒（Opgave 2 · Informationskløft）}
\small
两人一组面对面坐好，\textbf{严禁偷看对方题卡}，也\textbf{不可互相提示如何造句}。Kursist A 必须使用\textbf{第三人称（\textit{hun / Freja}）}对准下方两个主题的 \textbf{6 个信息格}各提 1 个聚焦问题（\textbf{切勿对卡片顶部的人物头像提问！}）。每提出 1 个语法与语义可接受的聚焦问题得 \textbf{1 分}，\textbf{满分 6 分，4--6 分及格}；本部分\textbf{仅提问方（A）计分}。
\end{tipbox}

\begin{tcolorbox}[enhanced, unbreakable, colback=NCDBgBlue, colframe=NCDNavy, coltitle=white, fonttitle=\bfseries, title={Kursist A Ark — Case: Freja（提问方专用 · 含图片线索与关键词）}, arc=3pt, boxrule=0.8pt, left=10pt, right=10pt, top=7pt, bottom=7pt]
\small
\centering \textbf{\color{NCDNavy}[卡片顶部人物肖像：Freja（请勿针对此头像提问）]}

\vspace{0.3em}
\raggedright
\renewcommand{\arraystretch}{1.3}
\begin{tabularx}{\linewidth}{@{} >{\raggedright\arraybackslash}X @{\hspace{0.5cm}} >{\raggedright\arraybackslash}X @{}}
\toprule
\textbf{\color{NCDNavy}Emne 1: Freja er i praktik på et hotel} & \textbf{\color{NCDNavy}Emne 2: Frejas sprogskole} \\
\midrule
\textbf{1. by / beliggenhed}\newline
\textcolor{NCDGray}{[图片提示：丹麦地图定位图标与海滨城市建筑]}\newline
\textbf{\color{NCDRed}?} &
\textbf{4. modul / niveau}\newline
\textcolor{NCDGray}{[图片提示：丹麦语教材封面与 Modul 1--6 阶梯图标]}\newline
\textbf{\color{NCDRed}?} \\
\midrule
\textbf{2. arbejdsopgaver}\newline
\textcolor{NCDGray}{[图片提示：酒店前台电脑、钥匙卡与早餐餐盘]}\newline
\textbf{\color{NCDRed}?} &
\textbf{5. dage / tid}\newline
\textcolor{NCDGray}{[图片提示：周历与晚间时钟图标]}\newline
\textbf{\color{NCDRed}?} \\
\midrule
\textbf{3. arbejdstider}\newline
\textcolor{NCDGray}{[图片提示：日历与工作时段钟表]}\newline
\textbf{\color{NCDRed}?} &
\textbf{6. transport til sprogskolen}\newline
\textcolor{NCDGray}{[图片提示：铁路与步行路线图标]}\newline
\textbf{\color{NCDRed}?} \\
\bottomrule
\end{tabularx}
\end{tcolorbox}

\begin{tcolorbox}[enhanced, unbreakable, colback=NCDBgGray, colframe=NCDGray, coltitle=white, fonttitle=\bfseries, title={Kursist B Ark — Case: Freja（回答方专用 · 事实信息卡片，请遮盖 A 卡作答）}, arc=2pt, boxrule=0.8pt, borderline west={3pt}{0pt}{NCDRed}, left=10pt, right=10pt, top=7pt, bottom=7pt]
\small
\renewcommand{\arraystretch}{1.3}
\begin{tabularx}{\linewidth}{@{} >{\raggedright\arraybackslash}X @{\hspace{0.5cm}} >{\raggedright\arraybackslash}X @{}}
\toprule
\textbf{\color{NCDTeal}Emne 1: Freja er i praktik på et hotel} & \textbf{\color{NCDTeal}Emne 2: Frejas sprogskole} \\
\midrule
\textbf{1.} Hotel Strandparken i \textbf{Helsingør}（靠近港口与克伦堡城堡）。 &
\textbf{4.} \textbf{Danskuddannelse 3, Modul 2}。 \\
\midrule
\textbf{2.} Hun tager imod gæster i receptionen, svarer på e-mails og hjælper med morgenmadsbuffeten. &
\textbf{5.} \textbf{Tirsdag og torsdag aften} kl. \textbf{17.30--20.00}。 \\
\midrule
\textbf{3.} \textbf{Mandag til fredag kl. 8.00--14.30}（周末不上班）。 &
\textbf{6.} Hun tager \textbf{toget (lokalbanen)} i 10 minutter og går 5 minutter fra stationen. \\
\bottomrule
\end{tabularx}
\end{tcolorbox}

% =====================================================================
% Del 4: Retteark, Detaljeret Analyse & Højscore-Modeller
% =====================================================================
\clearpage
\ncdsubsec{Del 4: Retteark, Detaljeret Analyse \& Højscore-Modeller}

\subsection*{4.1 官方 \textit{Retteark} 评分卡（Læsning 24 point + Mundtlig Opgave 2: 6 point）}

\begin{grammatikbox}{Læsning Rettenøgle — DU3 Modul 2 (Modultest 3)}
\small
\renewcommand{\arraystretch}{1.24}
\begin{tabularx}{\linewidth}{@{} >{\bfseries\color{NCDNavy}\raggedright\arraybackslash}p{2.3cm} >{\raggedright\arraybackslash}X >{\centering\arraybackslash}p{1.6cm} @{}}
\toprule
\textbf{大题} & \textbf{官方标准答案（Rettenøgle）} & \textbf{满分} \\
\midrule
Opgave 1 (NY) &
\textbf{1.} D \quad \textbf{2.} G \quad \textbf{3.} B \quad \textbf{4.} E \quad \textbf{5.} I \quad \textbf{6.} F \newline
（例题 \textbf{0.} A；多余干扰广告：\textbf{C} 与 \textbf{H}） & 6 point \\
\midrule
Opgave 2 &
须划掉的矛盾句（每句 1 分）：\newline
\textbf{1.} \textit{Ingen lagde mærke til den høje lyd.}\newline
\textbf{2.} \textit{Den store sorte hund gøede surt og bed alle gæsterne i benene.}\newline
\textbf{3.} \textit{Halsbåndet havde hverken navn eller adresse på.}\newline
\textbf{4.} \textit{Mikkel blev sur og nægtede at tage hunden med hjem.}\newline
\textbf{5.} \textit{Derfor aflyste Henrik festen og sendte alle vennerne hjem uden aftensmad.} & 5 point \\
\midrule
Opgave 3 &
\textbf{1.} og \quad \textbf{2.} retter \quad \textbf{3.} gave \quad \textbf{4.} sjov \quad \textbf{5.} hinanden \quad \textbf{6.} altid \newline
（例词：\textit{hedder}；多余干扰词：\textit{men, havde, når, mellem}） & 6 point \\
\midrule
Opgave 4 (NY) &
\textbf{1.} C \quad \textbf{2.} A \quad \textbf{3.} B \quad \textbf{4.} C \quad \textbf{5.} B \quad \textbf{6.} A \quad \textbf{7.} C \newline
（例题 \textbf{0.} B · 分布：A 两题、B 两题、C 三题） & 7 point \\
\midrule
\textbf{Læsning 总分} & \textbf{Bestået（及格）: 16--24 point} \quad | \quad \textbf{Ikke bestået（不及格）: 0--15 point} & \textbf{24 point} \\
\bottomrule
\end{tabularx}
\end{grammatikbox}

\begin{tipbox}{Mundtlig Opgave 2 (Informationskløft) 评分卡}
\small
\textbf{计分方式}：Kursist A 就卡片上 6 个栏目（\textit{1. by/beliggenhed, 2. arbejdsopgaver, 3. arbejdstider, 4. modul/niveau, 5. dage/tid, 6. transport til sprogskolen}）各发出 1 个第三人称（\textit{hun/Freja}）聚焦提问，每个有效问题得 \textbf{1 point}；满分 \textbf{6 point}，\textbf{4--6 point = Bestået}，0--3 point = Ikke bestået。Kursist B 在本题仅作答，不计分。
\end{tipbox}

% ---------------------------------------------------------------------
\subsection*{4.2 阅读 Opgave 1--4 逐题同义转述解析、全文中文精译与核心词表}

\begin{tipbox}{Opgave 1 解析：同义转述词对（Parafrasering）与干扰项拆解}
\small
\begin{itemize}
  \item \textbf{0 $\rightarrow$ A（例题）}：\textit{sortere bøger og hjælpe kunder}（整理书籍、服务顾客）$\Leftrightarrow$ \textit{sætte bøger på hylderne, tale med kunderne}；\textit{børn i nød} $\Leftrightarrow$ \textit{udsatte børn}（弱势儿童）。
  \item \textbf{1 $\rightarrow$ D}：\textit{om aftenen} + \textit{én ældre person en gang om ugen} + \textit{drikke kaffe og snakke} $\Leftrightarrow$ \textit{en eftermiddag eller aften hver uge} + \textit{en ældre dame eller herre} + \textit{over en kop kaffe}。排除 \textbf{H}（同样面向老人，但 H 是周三上午 10--12 点教 IT 手机操作，不符合埃米尔白天上班且想一对一喝咖啡聊天的需求）。
  \item \textbf{2 $\rightarrow$ G}：\textit{studerer matematik} + \textit{unge elever efter skoletid} + \textit{kun tid om tirsdagen} $\Leftrightarrow$ \textit{børn og unge 10--16 år} + \textit{hver tirsdag kl. 15.00--17.30} + \textit{forklare matematik}。
  \item \textbf{3 $\rightarrow$ B}：\textit{passe og lufte dyr i weekenden}（周末照看并遛动物）$\Leftrightarrow$ \textit{hundelufter ... lørdag eller søndag formiddag}。排除 \textbf{C}（虽然也在周末户外，但 C 是社区种菜除草，与动物无关）。
  \item \textbf{4 $\rightarrow$ E}：\textit{nyankommen flygtning} + \textit{øve dansk gennem almindelige samtaler} $\Leftrightarrow$ \textit{sprogmakker for en flygtning} + \textit{snakker sammen på dansk om arbejde, traditioner og livet i byen}。
  \item \textbf{5 $\rightarrow$ I}：\textit{ved vandet} + \textit{fjerne plastik og affald} $\Leftrightarrow$ \textit{på stranden} + \textit{samle plastik, dåser og andet affald ind, før det ender i havet}。排除 \textbf{C}（虽然同样户外环保，但仅限于社区花园种菜）。
  \item \textbf{6 $\rightarrow$ F}：\textit{torsdag formiddag} + \textit{på et sygehus} + \textit{byde velkommen ved indgangen og hjælpe med at finde afdeling} $\Leftrightarrow$ \textit{mandag til fredag kl. 8.30--13.00} + \textit{hospital} + \textit{ved hovedindgangen ... tage godt imod ... vise vej til den rigtige afdeling}。
  \item \textbf{多余干扰项}：\textbf{C}（社区园艺）与 \textbf{H}（文化馆老年 IT 辅导班）。
\end{itemize}
\end{tipbox}

\begin{transbox}
\small
\textbf{\color{NCDTeal}[Opgave 1 题目与广告中文精译]}\\
\textbf{人物需求 (0--6)}：
(0) 阿斯特丽德退休了，热爱小说与历史，想每周在温馨店铺用几个下午整理书籍、接待顾客，盈余捐给困境儿童。
(1) 埃米尔白天坐办公室，独居且晚上时间充裕，想念祖辈，希望每周固定与一位老人喝咖啡聊天。
(2) 娜迪娅在大学学数学，以前做过助教，想在放学后辅导青少年作业，但只有周二有空。
(3) 托比亚斯住市中心不准养宠物的小公寓，怀念乡下遛狗时光，想周末照看并遛动物。
(4) 萨米拉会丹麦语、阿拉伯语和英语，想在本地与新来的难民结对，通过日常聊天帮对方练丹麦语。
(5) 马兹关心海洋生态，整周坐电脑前，业余时间想去海边与大家一起清理塑料垃圾。
(6) 卡特琳娜学护理，周四上午有空，想去医院大门口迎接病患和家属并指引科室。\\
\textbf{招新启事 (A--I)}：
\textbf{[A]} 儿童二手书店招义工：售卖二手小说与童书，下午摆书、服务顾客与收银，收益归困境儿童。
\textbf{[B]} 动物收容所遛狗义工：招 18 岁以上义工于周六或周日上午带狗狗林间散步、喂水玩耍。
\textbf{[C]} 绿洲社区花园（干扰项）：周六上午一起浇水、除草、种有机蔬菜与建木箱并共进午餐。
\textbf{[D]} 老人陪访伙伴：每周一下午或晚间陪独居老人喝咖啡、散步或打牌约两小时。
\textbf{[E]} 难民语言伙伴：与新移民自由约定时间地点，用丹麦语聊工作、传统与城市生活。
\textbf{[F]} 医院导诊团队：周一至周五 8:30--13:00 在主入口迎接病患与访客并带路至相应科室。
\textbf{[G]} 青少年作业咖啡馆：每周二 15:00--17:30 在图书馆辅导 10--16 岁学生数学、物理或丹麦语作业。
\textbf{[H]} 文化馆老年 IT 辅导（干扰项）：每周三 10:00--12:00 帮助老人解决手机、网银和 MitID 问题。
\textbf{[I]} 洁净海滩行动：隔周日赴海滩捡拾塑料与罐头等垃圾，避免流入大海，现场提供手套、垃圾袋和热咖啡。
\end{transbox}

\begin{tipbox}{Opgave 2 解析：逻辑矛盾句逐段锁定}
\small
\begin{itemize}
  \item \textbf{(0) 例题}：7 月 14 日下午阳光明媚（\textit{Solen skinnede}），却说“整晚下大雪”（\textit{Det sneede kraftigt hele aftenen}），直接矛盾。
  \item \textbf{(1)}：上文刚说“大家都放下杯子好奇地朝草坪角落望去”，紧接着却说\textbf{“没人注意到那声巨响”}（\textit{Ingen lagde mærke til den høje lyd.}），自相矛盾。
  \item \textbf{(2)}：前文明确是“一只棕色小狗”（\textit{en lille brun hund}）且“高兴地摇尾巴舔手”，第 4 句却说\textbf{“那只黑色大狗生气地狂吠并咬了所有客人的腿”}，体貌与行为双双矛盾。
  \item \textbf{(3)}：前两句说路易丝在项圈上发现了“带名字和电话的金属牌，上面写着叫 Buster，住在 Bakkevej 12”，第 4 句却说\textbf{“项圈上既没有名字也没有地址”}，事实冲突。
  \item \textbf{(4)}：米克尔接到电话时“松了一口气”，到场后“微笑着解释”，小狗也奔向主人，第 4 句却说\textbf{“米克尔生气了并拒绝把狗带回家”}，人物态度完全相反。
  \item \textbf{(5)}：朋友们正围坐享用烤鸡沙拉、聊着趣事直到晚上十一点，第 4 句却说\textbf{“因此亨利取消了派对并让所有朋友没吃晚饭就回家了”}，与上下文完全矛盾。
\end{itemize}
\end{tipbox}

\begin{transbox}
\small
\textbf{\color{NCDTeal}[Opgave 2 中文精译（剔除矛盾句后的连贯译文）]}\\
(0) 7 月 14 日周六 17:30，亨利克和妻子路易丝邀请了六位好友来自家花园参加今夏第一次烧烤派对。蔚蓝的天空中阳光灿烂，客人们都已拿着冷饮坐在露台上。亨利克站在烤架旁煎香肠和鸡肉，夏日的美食香气飘过树篱。\\
(1) 六点半左右菜快好了，路易丝走进厨房去拿沙拉碗和新烤的面包。突然客人们听到花园门旁的高灌木丛后传来奇怪的声音。大家都放下杯子，好奇地望向草坪角落。就在这时，一只戴着红项圈的棕色小狗从灌木丛里跳出来，直奔烤架跑去。\\
(2) 小狗显然饿极了，急切地在亨利克腿边嗅来嗅去，还想够烤架上的香肠。亨利克小心翼翼地把小狗抱起来，以免它被滚烫的金属烫伤。它高兴地摇着尾巴，舔了舔亨利克的手。路易丝很快从厨房端出一碗冷水给这位不速之客。\\
(3) 路易丝仔细查看小狗的红项圈，发现了一块刻着名字和电话号码的小金属牌。上面写着小狗叫巴斯特，住在他们花园正后方的 Bakkevej 12 号。亨利克立刻拿出手机拨打了牌上的号码。一位名叫米克尔的年轻人接起电话，当亨利克告诉他巴斯特正安稳地坐在露台上时，他听起来如释重负。\\
(4) 五分钟后门铃响了，米克尔手里拿着一条长牵引绳站在门外。他微笑着解释说，刚才他在前院浇花时，巴斯特在树篱底下挖了个小洞钻了出来。亨利克领他到露台上，巴斯特立刻奔向自己的主人。客人们都友善地笑了起来，路易丝还在米克尔回家前请他喝了一杯冰柠檬水。\\
(5) 等米克尔和巴斯特走回 Bakkevej 后，朋友们终于在露台上入座开饭。烤鸡和沙拉美味极了，大家在晚餐期间津津有味地聊着刚刚的有趣插曲。气氛温馨又轻松，谁也不急着回家。直到十一点左右，客人们才道谢告别，在明亮的夏夜骑车回家。
\end{transbox}

\begin{tipbox}{Opgave 3 解析：句法搭配与干扰词排除}
\small
\begin{itemize}
  \item \textbf{(1) og}：固定并列结构 \textit{både ... og ...}（既……又……），对应前文 \textit{både julefrokost med sin familie, sine venner og sine kolleger}。排除 \textit{men}。
  \item \textbf{(2) retter}：复数限定词 \textit{mange} 后必须接可数名词复数“菜肴”（\textit{retter}），且后文列举了 \textit{sild, fiskefilet, frikadeller og flæskesteg}。
  \item \textbf{(3) gave}：不定冠词 \textit{en lille} 要求单数通性名词（\textit{en gave}，“礼物”），玩 \textit{pakkeleg}（抢礼物游戏）时每人带一件小礼物。
  \item \textbf{(4) sjov}：搭配 \textit{både \_\_\_ og højlydt} 需要并列形容词修饰通性主语 \textit{stemningen}（气氛很快变得既有趣又热闹）。
  \item \textbf{(5) hinanden}：高频固定搭配 \textit{at lære hinanden bedre at kende}（彼此更了解）。注意“互相了解”在丹麦语中必须用相互代词 \textit{hinanden}。
  \item \textbf{(6) altid}：情态动词 \textit{skal} + 主语 \textit{man} 之后紧接频率副词 \textit{altid}（总是）再接原形动词 \textit{huske}（如果玩到很晚，务必永远记得坐火车或打车回家，不要酒驾）。
  \item \textbf{4 个干扰词排除}：\textit{men}（无法配 \textit{både}）、\textit{havde}（过去时助动词，全文无此空位）、\textit{når}（从句连词，无相应位置）、\textit{mellem}（介词，无合法搭配）。
\end{itemize}
\end{tipbox}

\begin{transbox}
\small
\textbf{\color{NCDTeal}[Opgave 3 中文精译]}\\
十一月和十二月期间，许多丹麦人会举办一种叫做“圣诞午餐（julefrokost）”的特别聚会。人们既会和家人、朋友举办圣诞午餐，也会和工作中的同事们一起举办。在传统的圣诞午餐上，桌上会摆满许多菜肴，例如腌鲱鱼、炸鱼排、肉丸和脆皮烤猪肉。许多人配餐还会喝圣诞啤酒或烈酒（snaps），尽管有些人更喜欢喝水或汽水。饭后大家常常玩抢礼物游戏（pakkeleg），所有客人都事先带了一份小礼物。人们掷骰子赢取礼物，气氛很快变得既有趣又喧闹。公司圣诞午餐非常受欢迎，因为同事们可以穿着便装在办公室之外相聚。这是在新年到来之前彼此增进了解的好方式。如果狂欢到很晚，一定要记得坐火车或出租车回家，而不要开车。
\end{transbox}

\begin{tipbox}{Opgave 4 解析：三人学习方法定位与同义转述}
\small
\begin{itemize}
  \item \textbf{0 $\rightarrow$ B（例题）}：\textit{for nervøs til at tale dansk} $\Leftrightarrow$ \textit{meget genert ... bange for at sige noget forkert, så jeg turde næsten ikke åbne munden}。
  \item \textbf{1 $\rightarrow$ C (Elena)}：\textit{bor sammen med en dansker} $\Leftrightarrow$ \textit{flyttede ... sammen med min danske mand}。
  \item \textbf{2 $\rightarrow$ A (Marco)}：\textit{i pauserne på sit arbejde} $\Leftrightarrow$ \textit{I frokostpausen og over kaffen snakker vi kun dansk}。
  \item \textbf{3 $\rightarrow$ B (Yuna)}：\textit{litteratur for børn til at lære nye ord} $\Leftrightarrow$ \textit{låner tynde børnebøger ... de samme vigtige ord bliver gentaget ... nemmere kan huske dem}。
  \item \textbf{4 $\rightarrow$ C (Elena)}：\textit{øver sin skrivning hver eneste dag} $\Leftrightarrow$ \textit{Hver aften skriver jeg desuden fem korte sætninger i en notesbog}。
  \item \textbf{5 $\rightarrow$ B (Yuna)}：\textit{hjælpe ulønnet i en forretning}（在商店无偿帮忙）$\Leftrightarrow$ \textit{arbejder jeg frivilligt ... i en lokal genbrugsbutik}（\textit{ulønnet} = \textit{frivilligt}；\textit{forretning} = \textit{butik}）。
  \item \textbf{6 $\rightarrow$ A (Marco)}：\textit{ser fjernsyn med tekst på skærmen for at lære udtalen} $\Leftrightarrow$ \textit{ser danske krimiserier på tv med danske undertekster ... hvordan de skrevne ord lyder}。
  \item \textbf{7 $\rightarrow$ C (Elena)}：\textit{hører dansk tale, mens de laver husarbejde eller cykler} $\Leftrightarrow$ \textit{lytter til dansk radio og lette podcasts ..., mens jeg cykler, gør rent eller laver mad}。
\end{itemize}
\end{tipbox}

\begin{transbox}
\small
\textbf{\color{NCDTeal}[Opgave 4 三人自述中文精译]}\\
\textbf{A. 马尔科（Marco）}：两年前我从意大利来到奥胡斯担任工程师。第一年我几乎没学到什么丹麦语，尽管我每周去上两个晚上的语言学校，因为在工作开会时大家都讲英语。因此半年前我和两位丹麦同事做了个固定约定：午休和喝咖啡时我们彼此只说丹麦语；就算我说得很慢，他们也不能换成英语。晚上我还会看带丹麦语字幕的丹麦侦探剧，这极大地帮我理解书写的词在现实中怎么发音。我现在还是会犯语法错误，但终于敢和邻居以及在商店里说丹麦语了。\\
\textbf{B. 尤娜（Yuna）}：刚从韩国搬到欧登塞学设计时，我非常害羞。我怕说错话，几乎不敢用丹麦语开口。有一天老师建议我在校外找个练习的地方。现在我每周六在当地一家二手服装家具店做四小时义工。店里有很多健谈的顾客，通过关于价格和天气的简短对话，我对开口说话自信多了。此外，我还在图书馆借薄薄的儿童书回房间大声朗读。句子很简单，而且同样的常用词会反复出现，我记起来容易多了。\\
\textbf{C. 埃琳娜（Elena）}：一年半前我和我的丹麦丈夫从西班牙搬到罗斯基勒。我在家做平面设计师，服务西班牙客户，所以整天独自对着电脑，工作时听不到多少丹麦语。因此我每天在骑车、打扫卫生或做饭时，都会听近两小时的丹麦广播和简单播客。虽然并非每个词都听得懂，但我的耳朵已经习惯了语言的节奏和轻柔的词尾。此外，我每晚还在笔记本上写五个关于自己一天的短句，之后丈夫会帮我通读并讲解语序规则，让我不会忘记。这也让我在城里开口说丹麦语时从容多了。
\end{transbox}

\subsection*{核心备考词汇表（Gloser · 15 词精选）}

\begin{glosertable}
名/形 & en frivillig & \textit{(den frivillige, frivillige, de frivillige / frivillig, -t, -e)} & 志愿者；自愿的 \\
名词 & et overskud & \textit{(overskuddet, overskud, overskuddene)} & 盈余，利润；余力 \\
名词 & en borger & \textit{(borgeren, borgere, borgerne)} & 市民，居民 \\
动词 & at bekymre sig (om) & \textit{(bekymrer, bekymrede, har bekymret)} & 担心，关心 \\
动词 & at fjerne & \textit{(fjerner, fjernede, har fjernet)} & 清除，移走 \\
\midrule
名词 & en gæst & \textit{(gæsten, gæster, gæsterne)} & 客人，来宾 \\
名词 & en lyd & \textit{(lyden, lyde, lydene)} & 声音，响动 \\
名词 & et halsbånd & \textit{(halsbåndet, halsbånd, halsbåndene)} & （宠物）项圈 \\
动词 & at nægte & \textit{(nægter, nægtede, har nægtet)} & 拒绝 \\
名词 & en stemning & \textit{(stemningen, stemninger, stemningerne)} & 气氛，氛围 \\
\midrule
代词 & hinanden & \textit{(gensidigt stedord; kun flertal)} & 彼此，互相 \\
名词 & en aftale & \textit{(aftalen, aftaler, aftalerne)} & 约定，协议 \\
情态动词 & at turde & \textit{(tør, turde, har turdet)} & 敢，有勇气 \\
名词 & en ordstilling & \textit{(ordstillingen, ordstillinger, ordstillingerne)} & 语序，词序 \\
名词 & en ordrebekræftelse & \textit{(ordrebekræftelsen, -r, -rne)} & 订单确认函 \\
\end{glosertable}

% ---------------------------------------------------------------------
\subsection*{4.3 写作高分双档范文（Opgave 1 \& Bonus 1B）}

\begin{emailbox}{Opgave 1 · A2 稳过版（共 88 词｜覆盖全部 4 问，语序零失误）}
\textbf{Hej Jonas}

Tak for din mail! Det er rigtigt, at jeg er begyndt til en ny sport. Jeg går til badminton nu, og det er rigtig sjovt.

Jeg træner to gange om ugen, hver tirsdag og torsdag aften, i en hal tæt på mit hjem. Jeg valgte badminton, fordi jeg gerne ville have noget mere motion efter arbejde, og fordi min kollega anbefalede klubben.

De andre på holdet har taget godt imod mig. Alle er meget søde, og vi taler altid lidt sammen efter træningen.

\textbf{Mange hilsner}\\
\textbf{Lin}
\end{emailbox}

\begin{emailbox}{Opgave 1 · A2+ 高分示范版（共 118 词｜融汇本卷全部语法靶点）}
\textbf{Hej Jonas}

Tak for din besked! Ja, det passer faktisk, at jeg for tre uger siden begyndte til klatring i en lokal klub på Nørrebro. Jeg træner to aftener om ugen \da{for at} få mere energi i hverdagen.

Jeg valgte klatring, fordi min gode veninde Mette fortalte mig, \da{hvor spændende og udfordrende} \dabold{det er}. \da{Selvom} jeg i starten var lidt bange for højder, kan jeg allerede mærke, at min krop er blevet stærkere.

De andre på holdet har taget varmt imod mig. På holdet hjælper vi altid \dabold{hinanden} med at finde de rigtige greb, og træneren tager \da{sig} god tid til at vise \dabold{sin} teknik for os. Har du lyst til at prøve en dag?

\textbf{Mange hilsner}\\
\textbf{Lin}
\end{emailbox}

\begin{tipbox}{A2 vs A2+ 提分亮点点评}
\small
\begin{itemize}
  \item \textbf{词数安全线}：稳过版 88 词（高于 70 词底线 18 词），信息完整且无多余险招；高分版 118 词，叙事更有画面感。
  \item \textbf{目的与让步升级}：高分版自然使用 \da{for at + 原形} 表目的（代替零碎短句），并用前置让步从句 \da{Selvom jeg i starten var..., kan jeg allerede mærke...} 展现主句 V2 倒装功底。
  \item \textbf{间接引语与代词精准度}：嵌套间接疑问从句 \textit{fortalte mig, hvor spændende... det er}（严格不倒装！）；并同时准确区分相互代词 \dabold{hinanden}、反身代词 \da{sig} 与反身物主代词 \dabold{sin}（回指单数主语 \textit{træneren}）。
\end{itemize}
\end{tipbox}

\begin{emailbox}{Bonus Opgave 1B · 半正式客服维权范文（共 102 词｜全程复数 I / jer / jeres）}
\textbf{Hej Kundeservice}

Jeg skriver til \dabold{jer}, fordi jeg for 12 dage siden bestilte en vinterjakke i \dabold{jeres} netbutik (ordrenummer 33918). Ifølge ordrebekræftelsen skulle pakken blive leveret i løbet af 2--4 hverdage.

I onsdags fik jeg en besked om, at min ordre var pakket, men jeg har desværre ikke hørt noget fra \dabold{jer} siden i onsdags. \da{Selvom} jeg ved, at \dabold{I} har travlt, vil jeg gerne spørge, \da{hvornår} pakken bliver sendt, og \da{om} \dabold{I} kan sende mig et sporingsnummer. Jeg skal nemlig ud at rejse på fredag og har brug for jakken inden da.

På forhånd tak for hjælpen!

\textbf{Venlig hilsen}\\
\textbf{Lin Chen}
\end{emailbox}

\begin{tipbox}{Bonus 1B 易错避坑点}
\small
写信给公司、客服、房东机构或学校办公室时，丹麦语必须将收信方视为复数集体，全程统一使用 \textbf{\color{NCDRed}\textit{I / jer / jeres}}，\textbf{切忌中途滑回单数 \textit{du / dig / din}}！此外，问询配送时间与运单号时，\textit{hvornår pakken bliver sendt} 与 \textit{om I kan sende mig...} 均属间接问句，从句内绝对不可倒装。
\end{tipbox}

% ---------------------------------------------------------------------
\subsection*{4.4 口语高分逐字稿与信息差六问对照（Opgave 1 \& 2）}

\begin{tipbox}{Opgave 1 · 3A \textit{Mit studie}（131 词高分示范稿 + 重点语音 + Del 2 对答）}
\small
\textbf{Del 1 陈述示范（约 1 分钟 20 秒 · 131 词）：}\\
\textit{Jeg vil gerne fortælle om mit studie \ipa{ˈsduːˀdjə}. Jeg læser datalogi på Københavns Universitet \ipa{univɛɐ̯siˈtˢeːˀd} på Nørrebro, og jeg er på mit andet semester nu. Jeg har undervisning fra mandag til torsdag mellem klokken ni og tre, mens jeg bruger fredagen på at læse og lave opgaver hjemme.}

\textit{Mine vigtigste fag i år er programmering og matematik. Vi har dygtige og hjælpsomme \ipa{ˈjɛlbˌsʌmə} lærere, og vi taler altid meget afslappet med hinanden \ipa{hiˈnanən} i klassen.}

\textit{Jeg har gode studiekammerater \ipa{ˈsduːdjəˌkʰɑmɐˌʁɑːˀdɐ} fra både Danmark og udlandet. Vi har en lille studiegruppe, hvor vi spiser frokost i kantinen klokken tolv og drikker kaffe i pauserne for at snakke om vores projekter. Hver morgen cykler jeg tyve minutter langs søerne hen til universitetet, selvom \ipa{ˈsɛlˌʌmˀ} jeg selvfølgelig tager metroen, når det regner kraftigt. Jeg er rigtig glad for mit studie.}

\vspace{0.3em}
\textbf{Del 2 考官追问与高分对答：}
\begin{enumerate}
  \item \textbf{Eksaminator:} \textit{Hvilket fag kan du bedst lide, og hvorfor?}\\
  \textbf{Kursist:} \textit{Jeg kan bedst lide programmering, fordi vi laver praktiske projekter i grupper, og det er spændende at løse problemer sammen.}
  \item \textbf{Eksaminator:} \textit{Hvilket sprog taler du med dine studiekammerater i pauserne?}\\
  \textbf{Kursist:} \textit{Til undervisningen taler vi ofte engelsk, men i pauserne prøver jeg at tale dansk med mine danske studiekammerater, så jeg kan øve mig.}
  \item \textbf{Eksaminator:} \textit{Hvad vil du gerne arbejde med, når du er færdig?}\\
  \textbf{Kursist:} \textit{Når jeg er færdig, vil jeg gerne arbejde som softwareudvikler i en dansk virksomhed, hvor vi bruger både dansk og engelsk.}
  \item \textbf{Eksaminator:} \textit{Er det svært at have tid til både studie og dansk sprogskole?}\\
  \textbf{Kursist:} \textit{Ja, nogle gange har jeg travlt, men fordi jeg har fri om fredagen, kan jeg planlægge min tid godt.}
\end{enumerate}
\end{tipbox}

\begin{tipbox}{Opgave 1 · 3B \textit{Min fritidsaktivitet: Vinterbadning}（132 词高分示范稿 + 重点语音 + Del 2 对答）}
\small
\textbf{Del 1 陈述示范（约 1 分钟 20 秒 · 132 词）：}\\
\textit{Jeg vil gerne fortælle om min fritidsaktivitet, som er vinterbadning \ipa{ˈvenˀdɐˌbæðneŋ}. Jeg er medlem af vinterbadeklubben Isfuglen ved Amager Strandpark i København. Jeg kommer der hver tirsdag og torsdag morgen fra klokken syv til kvart i otte før arbejde, og nogle gange også søndag formiddag.}

\textit{Det koster seks hundrede og halvtreds \ipa{halˈtˢʁɛs} kroner om året, og prisen giver adgang til både omklædning og en varm sauna \ipa{ˈsɑʊ̯na}. Jeg bader næsten altid sammen med to søde naboer og mange andre faste medlemmer, så det er et rigtig hyggeligt fællesskab \ipa{ˈfɛləsˌsgæːˀb}. Først hopper vi kort i det kolde vand, bagefter sidder vi femten minutter i saunaen, og til sidst drikker vi varm te og snakker sammen.}

\textit{Jeg cykler derhen på tolv minutter med mit badetøj og et stort håndklæde \ipa{ˈhʌnˌkʰlɛːðə} i rygsækken \ipa{ˈʁøɡˌsɛɡ}. Det giver mig masser af ny energi.}

\vspace{0.3em}
\textbf{Del 2 考官追问与高分对答：}
\begin{enumerate}
  \item \textbf{Eksaminator:} \textit{Hvorfor begyndte du at vinterbade? Var det ikke meget koldt i starten?}\\
  \textbf{Kursist:} \textit{Min nabo spurgte mig, om jeg ville prøve en dag i oktober. I starten var vandet et chok, men nu har min krop vænnet sig til det.}
  \item \textbf{Eksaminator:} \textit{Hvordan føles kroppen efter saunaen og det kolde vand?}\\
  \textbf{Kursist:} \textit{Man vågner fuldstændig op og bliver i rigtig godt humør hele dagen.}
  \item \textbf{Eksaminator:} \textit{Hvad siger dine venner eller din familie til, at du bader om vinteren?}\\
  \textbf{Kursist:} \textit{Min familie i mit hjemland synes, det lyder lidt skørt, men mine danske venner synes, det er helt normalt!}
  \item \textbf{Eksaminator:} \textit{Laver I også noget andet socialt i klubben ud over at bade?}\\
  \textbf{Kursist:} \textit{Ja, en gang om måneden spiser vi morgenmad sammen om søndagen efter saunaen.}
\end{enumerate}
\end{tipbox}

\begin{grammatikbox}{Opgave 2 · \textit{Freja} 信息差六问全解（直接问句 vs 礼貌间接问句 + B 卡答案）}
\small
\renewcommand{\arraystretch}{1.28}
\begin{tabularx}{\linewidth}{@{} >{\bfseries\raggedright\arraybackslash}p{2.9cm} >{\raggedright\arraybackslash}p{3.8cm} >{\raggedright\arraybackslash}p{4.3cm} >{\raggedright\arraybackslash}X @{}}
\toprule
\textbf{提示词格位} & \textbf{直接问法（稳过 1 分）} & \textbf{礼貌间接问法（秀语序 1 分）} & \textbf{Kursist B 参考回答} \\
\midrule
1. by / beliggenhed &
\textit{Hvor ligger det hotel, hvor Freja er i praktik?} &
\textit{Kan du fortælle mig, \da{hvor} \dabold{Frejas praktiksted ligger}?} &
\textit{Hun er i praktik på Hotel Strandparken i Helsingør.} \\
\midrule
2. arbejdsopgaver &
\textit{Hvad er hendes arbejdsopgaver på hotellet?} &
\textit{Jeg vil gerne vide, \da{hvad} \dabold{hun skal lave} på hotellet.} &
\textit{Hun tager imod gæster i receptionen, svarer på e-mails og hjælper med morgenmaden.} \\
\midrule
3. arbejdstider &
\textit{Hvornår arbejder hun, og arbejder hun i weekenden?} &
\textit{Ved du, \da{hvilke dage og tider} \dabold{hun arbejder}?} &
\textit{Hun arbejder mandag til fredag fra kl. 8.00 til 14.30.} \\
\midrule
4. modul / niveau &
\textit{Hvilket modul går hun på på sprogskolen?} &
\textit{Kan du sige mig, \da{hvilket modul} \dabold{hun går på}?} &
\textit{Hun går på Danskuddannelse 3, Modul 2.} \\
\midrule
5. dage / tid &
\textit{Hvornår går hun til dansk på sprogskolen?} &
\textit{Jeg vil gerne vide, \da{hvornår} \dabold{hun har undervisning} på sprogskolen.} &
\textit{Hun går på sprogskole tirsdag og torsdag aften kl. 17.30--20.00.} \\
\midrule
6. transport &
\textit{Hvordan kommer hun hen til sprogskolen?} &
\textit{Kan du fortælle mig, \da{hvordan} \dabold{hun kommer} til sprogskolen?} &
\textit{Hun tager toget i 10 minutter og går 5 minutter fra stationen.} \\
\bottomrule
\end{tabularx}

\vspace{0.4em}
\textbf{\color{NCDNavy}第三人称信息转述汇总（Freja · 84 ord）：}\\
\textit{Freja er i praktik på Hotel Strandparken i Helsingør, hvor hun arbejder fra mandag til fredag mellem klokken 8.00 og 14.30. Hun tager imod gæster i receptionen, svarer på e-mails og hjælper med morgenmadsbuffeten. Ved siden af sin praktik går Freja på sprogskole på Danskuddannelse 3, Modul 2. Hun har undervisning tirsdag og torsdag aften fra klokken 17.30 til 20.00. Når hun skal til sprogskole, tager hun toget i ti minutter og går fem minutter fra stationen. Selvom ugen er travl, trives hun godt.}
\end{grammatikbox}
